\section{L'électrification et le substrat commun (~1880--1930)}\label{sec:electrification}
% =============================================================================

% §1.7 — Ouverture avec analyse des mécanismes structurels
% DRAFT v3 — 22 avril 2026
% Cinq mécanismes : (1) information comme condition de viabilité du système de puissance,
% (2) load factor comme problème informationnel, (3) infrastructure partagée et path dependence,
% (4) monopole naturel par rendements d'échelle, (5) reverse salient.
% La bifurcation power/information n'est PAS posée ici — elle émerge dans le bilan.
% David/Devine → §1.8.
% Hughes cité via secondaires + VdK ; \wip{} là où le primaire est indispensable.

Les cas qui précèdent partageaient un trait commun que le récit n'avait pas eu besoin de nommer~: aucun d'entre eux ne dépendait du réseau électrique pour fonctionner. Le télégraphe de Morse tirait son énergie de piles voltaïques~; les calculatrices de Thomas de Colmar et de Burroughs étaient actionnées à la main~; les tabulatrices Hollerith consommaient un courant dont le coût représentait 1,5\,\% du budget d'exploitation. L'électricité traversait ces cas comme un détail technique, pas comme une condition de fonctionnement. La période qui s'ouvre ici change la donne~: la radiodiffusion, l'usage massif des tabulatrices dans les grandes organisations, et bientôt le calcul électronique supposent tous l'existence d'un réseau électrique fiable, centralisé et continu. Le réseau électrique n'est pas une technique de l'information au sens de ce chapitre~; mais la manière dont il s'est constitué, entre 1880 et 1930, révèle des mécanismes qui intéressent directement l'analyse, parce qu'ils mettent en jeu, pour la première fois, l'articulation entre des fonctions informationnelles et un système de puissance au sein d'un même artefact technique.

\dimmark{D3/D7}{+}%
Le 4~septembre 1882, la Pearl Street Central Generation Station ouvrit dans le sud de Manhattan. Ce que Thomas Edison avait conçu n'était pas une lampe à incandescence~; d'autres l'avaient précédé, de Swan en Angleterre à Sawyer aux États-Unis. C'était un système complet, calqué sur l'architecture de l'industrie du gaz~: une station centrale produisant l'électricité, un réseau de câbles souterrains la distribuant, un compteur kilowattheure la mesurant, un tarif à l'usage la facturant \parencite{kooij_electric_light_2015}. \textcite{hughes_networks_1983} a montré que cette conception systémique était la contribution décisive d'Edison, celle qui le distinguait de ses prédécesseurs~: non pas l'invention d'un composant, mais l'architecture d'un ensemble dont chaque élément était dimensionné par rapport aux autres\footnote{Hughes qualifie Edison de \emph{system builder} et organise l'évolution du système électrique en quatre phases~: invention et développement, transfert technologique, croissance du système (avec les \emph{reverse salients}), et \emph{technological momentum}. Le modèle est appliqué systématiquement à trois trajectoires nationales~: New York, Berlin et Londres \parencite{hughes_networks_1983}.}. Six générateurs de vingt-sept tonnes chacun, alimentés par des chaudières consommant cinq tonnes de charbon et 43\,500~litres d'eau par jour, fournissaient 600~kilowatts de courant continu à 110~volts. Quatre cents lampes brûlèrent le premier jour chez quatre-vingt-cinq clients, dont Drexel, Morgan \& Company et le \emph{New York Times}~; un an plus tard, 8\,573~lampes desservaient 513~clients \parencite{kooij_electric_light_2015}.

\dimmark{D6/D7}{+}%
Or le système de Pearl Street ne pouvait pas fonctionner sans un dispositif informationnel qui lui était constitutif. Avant le compteur kilowattheure (brevet US n\textsuperscript{o}~242,901, 14~juin 1881), Edison facturait au forfait par lampe installée~; le client payait le même montant qu'il consommât beaucoup ou peu, ce qui rendait impossible toute incitation à la sobriété et toute adaptation de la capacité à la demande réelle. Le compteur transformait un flux physique continu en un chiffre discret, lisible, facturable~: c'est un dispositif de traitement de l'information, logé au c\oe{}ur d'un système de puissance, sans lequel le modèle économique de la centrale n'existe pas. La chaîne causale est la suivante~: le compteur rend possible le tarif à l'usage~; le tarif à l'usage rend possible la diversification de la base de clients~; la diversification rend possible l'optimisation du taux d'utilisation de la capacité installée. C'est la première fois dans le parcours du chapitre que l'information n'agit ni sur un marché extérieur (comme le télégraphe créant le marché du coton chez \textcite{steinwender_real_2018}) ni sur la coordination interne d'une organisation (comme les six formes de McCallum sur l'Erie Railroad)~: elle rend un système de puissance \emph{viable}.

\dimmark{D2}{+}%
Le système d'Edison souffrait cependant d'une limitation physique qui allait déterminer la suite. Le courant continu ne pouvait être transporté qu'à basse tension, ce qui imposait la proximité entre la centrale et les clients~: Pearl Street ne desservait qu'un rayon d'environ un mile\footnote{La perte en ligne croît avec le carré du courant et la longueur du conducteur. En courant continu à 110~volts, doubler la distance de desserte exigeait de quadrupler la section des câbles de cuivre, rendant le coût prohibitif au-delà d'un mile environ.}. \textcite{hughes_networks_1983} nomme ce type de contrainte un \emph{reverse salient}~: le composant du système dont l'insuffisance empêche l'ensemble de progresser, par analogie avec le segment d'un front militaire qui n'avance pas au même rythme que le reste. La distance de transmission était le \emph{reverse salient} du système Edison~; sa résolution vint de l'extérieur du système, avec l'adoption du courant alternatif portée par Westinghouse et Tesla dans les années~1890. Le transformateur permettait de transporter l'électricité à haute tension sur de longues distances et de la distribuer à basse tension chez le client~: le \emph{reverse salient} était levé, mais au prix d'un remplacement complet de l'architecture du réseau\footnote{La \emph{War of Currents} entre Edison (DC) et Westinghouse (AC), qui opposa les deux systèmes pendant une décennie, est un cas où le choix technique n'avait pas de solution unique~: \textcite{hughes_networks_1983} montre que le système DC aurait pu survivre avec des améliorations incrémentales (le système à trois fils, les batteries de stockage), et que la victoire de l'AC ne fut ni inévitable ni purement technique. \textcite{granovetter_making_1998} soutiennent que la structure de l'industrie électrique fut construite socialement, pas déterminée techniquement.}. Le concept de \emph{reverse salient} est utile au-delà du cas Edison~: il sera repris en fin de chapitre pour formuler l'hypothèse que l'énergie est, au vingt et unième siècle, le \emph{reverse salient} du système informationnel lui-même.

\dimmark{D2/D10}{+}%
Samuel Insull, ancien secrétaire d'Edison devenu directeur de la \emph{Chicago Edison Company}, comprit que la clé du réseau n'était pas la lampe mais le \emph{load factor}~: le taux d'utilisation de la capacité installée. Une centrale qui ne tourne qu'aux heures de pointe produit de l'électricité coûteuse~; une centrale qui tourne en continu produit de l'électricité bon marché. Le problème était de nature informationnelle~: il fallait savoir qui consommait quand, comment la demande se répartissait dans la journée, quelles catégories de clients pouvaient lisser la charge. En diversifiant la base (résidentiel le soir, industriel le jour, tramways en continu), Insull fit baisser le coût unitaire~; en négociant des tarifs dégressifs, il augmenta la consommation et justifia des centrales de plus en plus grandes\footnote{\wip{Développer avec Hughes 1983, ch.~8 (\emph{Chicago~: The Dominance of Technology})~: le modèle Insull, les rendements d'échelle, le monopole naturel régulé par les \emph{public utility commissions} créées à partir de 1907. La comparaison avec Berlin (ch.~7, Rathenau, coordination technologie-politique) et Londres (ch.~9, primauté du politique) produirait trois trajectoires d'électrification, comme les sept trajectoires du télégraphe en §1.2, convergeant vers un résultat commun~: le système intégré à grande échelle.}}. Le mécanisme de monopole qui en résultait était distinct des deux précédents~: ni les rendements de réseau du télégraphe, ni le verrouillage par le consommable d'IBM, mais les rendements d'échelle dans la génération, combinés à une régulation publique qui garantissait au monopoleur son territoire en échange d'un contrôle des tarifs. \textcite{hughes_networks_1983} nomme \emph{technological momentum} la résistance au changement d'un système dans lequel des investissements massifs, des institutions éducatives, des sociétés professionnelles et des cadres réglementaires se sont co-construits autour d'une architecture technique~: une fois le réseau en place, ses composants co-évoluent et résistent au remplacement. Le réseau électrique acquit ce momentum avant 1914~; les techniques de l'information, elles, allaient désormais évoluer dans le milieu qu'il avait créé.

\dimmark{D10/0-TER}{+}%
Un dernier trait mérite d'être relevé, parce qu'il n'est visible qu'avec le recul que donne la traversée des cas précédents. Edison n'était pas électricien de formation~: il était opérateur de télégraphe. Sa carrière commença par le multiplexage télégraphique~; le quadruplex (1874), qui permit de transmettre quatre messages simultanément sur un seul fil, fit sa réputation et sa fortune initiale \parencite{kooij_electric_light_2015}. C'est la maîtrise du circuit électrique, acquise dans le domaine de la communication, qu'il transposa dans le domaine de la puissance. Les poteaux du télégraphe portèrent les premiers fils de distribution électrique. Western Electric, qui fabriquait du matériel télégraphique, fabriqua ensuite du matériel téléphonique, puis du matériel électrique. Les mêmes chaînes d'approvisionnement en cuivre servaient les deux réseaux. Cette parenté n'est pas une coïncidence~: elle signale que les deux trajectoires de l'électricité~--- la puissance et l'information~--- naissent dans les mêmes ateliers, avec les mêmes matériaux, par les mêmes ingénieurs, avant de diverger institutionnellement en industries séparées\footnote{La chronologie renforce l'observation~: le télégraphe électrique (1837) précède la dynamo auto-excitatrice (1866) de trente ans. La trajectoire de l'information précède celle de la puissance. Cette antériorité et cette parenté seront reprises dans le bilan de cette section.}. Au tournant du siècle, la divergence était consommée~: l'industrie électrique et l'industrie des télécommunications avaient des acteurs, des régulateurs et des revues savantes distincts. Mais le substrat restait commun, et la dépendance des techniques de l'information à ce substrat allait s'intensifier dans chacun des cas qui suivent.

\dimmark{0-TER}{+}%
Cette intensification restera pourtant invisible pendant quatre décennies supplémentaires, grâce à un mécanisme que \textcite{dennard_design_1974} ont formalisé sans en mesurer les conséquences énergétiques. Le \emph{MOSFET scaling} de Dennard établit que si les dimensions d'un transistor diminuent d'un facteur $k$, la tension et le courant diminuent du même facteur, de sorte que la densité de puissance~--- watts par centimètre carré~--- reste constante. En pratique, le voltage des transistors passe de 5~volts au début des années~1990 à 1,2~volt vers~2005 \parencite{roussilhe_phase_2025}. Tant que cette loi tient, la puissance de calcul peut croître exponentiellement (selon la trajectoire de \textcite{moore_cramming_1965}) sans que la consommation électrique des composants n'augmente proportionnellement. C'est le bouclier qui permet au discours de dématérialisation de se maintenir~: l'ordinateur est perçu comme de plus en plus puissant et de moins en moins énergétivore par opération, ce qui est vrai~--- mais la consommation agrégée croît parce que le volume de calcul croît plus vite que l'efficacité. Vers~2005, Dennard atteint sa limite physique aux alentours de 3~GHz~; l'industrie contourne l'obstacle par la multiplication des c\oe{}urs (\emph{multi-core}), mais la digue électrique a cédé. La puissance maximale (TDP) des processeurs, stabilisée entre 75 et 150~watts pendant la décennie précédente, recommence à croître~; avec les GPU dédiés à l'IA générative, elle atteint 700~watts en~2024 et les feuilles de route projettent 1\,400~watts d'ici~2026 \parencite{roussilhe_phase_2025, sun_chip_2024}. Le \emph{scaling} de Dennard aura ainsi joué le rôle d'un facteur d'occultation au sens de \textcite{ensmenger_environmental_2018}~: pendant quarante ans, il a maintenu le substrat électrique hors du champ de vision des économistes de l'ICT, non par dissimulation délibérée mais parce que le mécanisme physique rendait la question énergétique littéralement sans objet à l'échelle du composant. La section~\ref{sec:computing} reprend là où ce bouclier cède.

\subsection{Berlin, Londres et la divergence des trajectoires}\label{subsec:trajectoires_europe}

\dimmark{D2/0-GOV}{+}%
Le système Edison que la section précédente a décrit ne s'est pas déployé à l'identique dans les trois capitales où il fut transféré. \textcite{hughes_networks_1983} a documenté les trajectoires divergentes de New York, Berlin et Londres entre 1882 et 1914~; leur comparaison révèle que la forme institutionnelle, et non la supériorité technique, détermine le rythme d'électrification. À Berlin, Emil Rathenau, qui avait acquis les droits Edison pour l'Allemagne après avoir visité l'exposition de Paris en 1881, fonde la \emph{Deutsche Edison Gesellschaft} (1883), devenue \emph{Allgemeine Elektricitäts-Gesellschaft} (AEG) en 1887. La concession municipale de 1884 lui accorde le monopole de la distribution dans un périmètre défini, en échange du partage des revenus avec la ville~; AEG contrôle simultanément la fabrication des équipements et l'exploitation du réseau, une intégration verticale que ni Edison à New York ni aucun concurrent à Londres n'atteindra\footnote{Siemens \& Halske, le concurrent berlinois d'AEG, adopte une stratégie complémentaire plutôt que frontale~: les deux firmes se partagent le marché allemand et coopèrent dans les projets d'infrastructure à grande échelle. \textcite{hughes_networks_1983}, ch.~7, documente cette coordination comme un trait distinctif du modèle allemand.}. En 1914, la Berliner Elektrizitäts-Werke (BEW), filiale d'AEG, dessert un rayon de trente kilomètres avec six centrales et une capacité installée de 155\,000~kW~; le taux de charge (\emph{load factor}) y est supérieur à celui de Chicago et bien au-dessus de celui de Londres.

\dimmark{0-GOV}{+}%
À Londres, la trajectoire fut inverse. L'\emph{Electric Lighting Act} de 1882, promu par Joseph Chamberlain au Board of Trade, imposa aux compagnies privées une clause de rachat municipal à vingt et un ans au prix des actifs isolés (\emph{scrap iron clause})~: la municipalité pouvait racheter câbles, dynamos et bâtiments à leur valeur séparée, sans prime pour la valeur du système en fonctionnement\footnote{L'expression \emph{scrap iron} provient des débats parlementaires de 1882~: le secrétaire du Board of Trade expliqua que le rachat se ferait à la valeur des «~briques comme briques et du mortier comme mortier~» (\emph{bricks as bricks and mortar as mortar}). Voir \textcite{hughes_networks_1983}, p.~57.}. Cette clause gela l'investissement privé~: pourquoi immobiliser du capital dans une infrastructure dont la valeur serait amputée au moment du rachat~? En 1883, soixante-neuf autorisations provisoires furent accordées~; en 1884, quatre~; en 1885, aucune. La station Edison de Holborn Viaduct, ouverte en 1882, fut abandonnée par Edison \& Swan en 1886. Quand le Parlement amenda la loi en 1888 pour porter la durée de concession à quarante-deux ans, il était trop tard~: les villes avaient acquis le droit de gérer elles-mêmes l'éclairage, et la fragmentation s'enracina. En 1918, Londres comptait encore soixante-cinq compagnies d'électricité, dix fréquences différentes et vingt-quatre tensions~; l'unification ne viendrait qu'avec le \emph{Central Electricity Board} de 1926.

\dimmark{D2/D10}{*}%
La comparaison Hughes établit un résultat qui dépasse le cas de l'électricité. La trajectoire technique du réseau n'explique pas sa forme institutionnelle~: Berlin et Londres partaient du même artefact (le système Edison), disposaient d'ingénieurs compétents et de capitaux disponibles~; c'est la configuration politique qui discrimina les trajectoires. À Berlin, la coalition Rathenau--banques d'investissement--municipalité avait le pouvoir d'imposer une concession stable~; à Londres, la coalition aristocratie--science--industrie se heurta à des intérêts contradictoires au Parlement. Hughes résume~: \og~The most penetrating explanation for the failure in London and the success in Berlin is neither technological nor economic; it is political.\fg{}\footnote{\textcite{hughes_networks_1983}, p.~78.} Ce résultat est exactement parallèle à celui que la section~\ref{sec:telegraph} avait établi pour le télégraphe~: sept trajectoires nationales convergeant vers le monopole par sept mécanismes différents, dont la variété tient aux configurations politiques locales et non aux propriétés techniques du réseau. La dimension \textbf{0-GOV} n'est pas un résidu~; elle est, dans les deux cas, le facteur déterminant de la forme institutionnelle.

\emergence{La divergence Berlin/Londres confirme que les réseaux électriques, comme les réseaux télégraphiques, convergent vers des formes institutionnelles comparables (monopole régulé, intégration verticale, contrôle étatique) par des trajectoires politiquement déterminées. La supériorité technique n'explique ni le succès de Berlin ni l'échec de Londres~: c'est la capacité des coalitions d'acteurs à stabiliser un cadre institutionnel favorable à l'investissement long qui discrimine les trajectoires. Le mécanisme vaut pour 0-TER~: là où l'électrification est retardée par la fragmentation, les techniques de l'information restent plus longtemps sur substrat mécanique ou manuel, ce qui retarde aussi leur couplage au réseau de puissance.}


% --- Les subsections continuent ici : phonographe, IBM tabulatrices, Shannon, Enigma ---


\subsection{Le phonographe et les supports de mémoire}\label{subsec:phonographe}

\dimmark{0-TER}{*}%
Edison invente le phonographe en 1877, un an après le téléphone de Bell. Il est frappant que le même homme ait conçu, en l'espace d'une décennie, le compteur kilowattheure (information \emph{pour} la puissance), le quadruplex (information \emph{pour} la transmission) et le phonographe (information \emph{pour} la mémoire)~: trois fonctions en un seul atelier, non comme programme délibéré mais comme convergence d'un même génie du circuit électrique vers ses trois usages possibles. Le gramophone de Berliner (1887) et le Télégraphone de Poulsen (1898) prolongent la même logique~: des informations peuvent être fixées sur un support matériel et relues à volonté, sans consommer de courant entre les deux actes.

Ce trait est précisément ce qui distingue le régime pré-numérique du stockage de tous les régimes qui suivront. La cire, le disque de gomme-laque, le fil magnétique, le film photographique retiennent l'information sans énergie~: l'inscription est une déformation physique permanente du support, et la lecture en est la traduction mécanique ou optique. Le profil énergétique de STORE dans toute cette période est celui d'une ressource à constitution ponctuelle et à mobilisation ponctuelle~; entre les deux, l'énergie est nulle. Cette passivité n'est pas une limitation technique en attente de correction~: c'est une propriété structurelle des supports analogiques, qui tient à leur nature physique. Elle cessera d'être la norme avec la mémoire à accès dynamique (DRAM, Intel 1103, 1970), dont les cellules doivent être rafraîchies des centaines de fois par seconde sous peine de perdre leur contenu. La frontière entre le régime passif et le régime actif-permanent de STORE correspond à cette transition, qui sera traitée dans la section suivante.


\subsection{La radio et la télévision de masse}\label{subsec:broadcasting_mass}

\dimmark{D1/D5}{+}%
La radiodiffusion de masse ~--- NBC en 1926, BBC en 1927, puis les systèmes nationaux européens tout au long des années~1930~--- accomplit quelque chose que le téléphone, le télégraphe et même la presse n'avaient pas réalisé sous cette forme~: la transmission simultanée d'un même signal à des millions de récepteurs sans que le nombre de récepteurs affecte ni la qualité du signal ni son coût de production. C'est la première réalisation technologique d'un bien strictement non rival~: la énième auditrice d'un concert radiodiffusé ne diminue pas ce que reçoit la première. Les ventes de récepteurs radio aux États-Unis illustrent l'ampleur du phénomène~: soixante millions de dollars en 1922, huit cent cinquante millions en 1929, en l'espace de sept ans seulement.

\dimmark{D8}{+}%
Cette croissance pose immédiatement la question du modèle économique. Un signal radiodiffusé ne peut pas être rendu excluable par un mécanisme technique simple~: quiconque possède un récepteur peut capter l'émission sans payer. Deux solutions émergent~: la redevance, adoptée en Europe dans le cadre d'opérateurs publics (BBC, ORTF, ARD), et la publicité, choisie aux États-Unis par les réseaux commerciaux. Ces deux solutions impliquent deux régimes de propriété et deux rapports à la demande~: la redevance découple la production du contenu du consentement à payer des auditeurs individuels~; la publicité découple le coût du signal de sa valeur pour l'auditeur, au profit d'un tiers payeur dont l'intérêt n'est pas le programme mais l'attention. Le rapport ICT-énergie dans les deux cas est identique~: le réseau électrique est nécessaire à la réception domestique de masse, mais ce coût est diffus et invisible~--- il est noyé dans la facture d'électricité du ménage, pas dans le coût d'exploitation du diffuseur.


\subsection{IBM et les tabulatrices dans les grandes organisations}\label{subsec:ibm_tab}

\dimmark{D2/D5}{*}%
La période 1924--1935 voit se refermer le cycle ouvert par Hollerith~: la «~Computing-Tabulating-Recording Company~», rebaptisée IBM («~International Business Machines~») par Thomas Watson en 1924, consolide une position de quasi-monopole que ni la concurrence de Remington Rand ni l'antitrust de 1936 ne parviendront à défaire. Deux événements accélèrent ce mouvement. Le premier est technique~: en 1928, IBM introduit la carte à quatre-vingt colonnes à trous rectangulaires, brevetée, qui brise l'interopérabilité entre les systèmes Hollerith et Powers et verrouille les clients existants sur les matériels IBM. Le second est institutionnel~: en 1935, le Social Security Act impose à l'administration américaine de tenir les comptes individuels de cotisation de plus de vingt-six millions de travailleurs, un volume de données qui excède la capacité de tout système manuel~; la tabulatrice, dont IBM est le fournisseur dominant, devient dès lors une infrastructure d'État, financée par les impôts et déployée à une échelle que le marché privé n'aurait pas atteinte aussi rapidement. Cette articulation entre crise administrative de la grande organisation et déploiement de la mécanographie comme infrastructure de coordination s'inscrit dans ce que \textcite{beniger_control_1986} a appelé la \emph{control revolution}~: la révolution informationnelle n'y serait pas une rupture post-industrielle mais une réponse continue à la crise de contrôle déclenchée par la révolution industrielle elle-même, dont les chemins de fer du XIX\textsuperscript{e}~siècle, l'organisation chandlérienne et la mécanographie hollerithienne constituent les phases successives. Le cadre est retenu ici comme lecture disponible~; la décomposition par fonctions TRANSMIT/STORE/COMPUTE et le suivi du couplage énergétique~0-TER en restent les apports propres au présent chapitre.

\dimmark{D4}{+}%
Le cas du \emph{Nautical Almanac Office} que étudie \textcite{daston_calculations_2018} éclaire ce que la mécanisation fait réellement au travail de calcul. Entre 1928 et 1937, le bureau mécanise progressivement le calcul des éphémérides~: le budget annuel passe de cinq cents livres sterling à six cent sept, soit une hausse de vingt et un pour cent, tandis que le nombre de collaborateurs augmente de même. La tabulatrice ne remplace pas les calculateurs humains~: elle déplace le goulot d'étranglement. Ce qui manquait avant la mécanisation était le calcul arithmétique brut~; ce qui manque après est ce que les superviseurs de l'époque appellent \og~analytical intelligence~\fg{}~: la capacité à formuler les problèmes, à vérifier les résultats, à concevoir les séquences d'opérations. Vingt à trente pour cent du temps de travail est consacré à la supervision et au contrôle~; la mécanisation déplace la pénurie vers ces tâches. Le mécanisme est déjà celui qu'\textcite{autor_skills_2003} formalisera soixante ans plus tard~: la mécanisation des tâches routinières codifiables augmente la demande de tâches cognitives non routinières, elle ne la supprime pas.

\dimmark{0-TER}{*}%
Le profil énergétique de ce régime reste négligeable. Les tabulatrices IBM des années~1930 consomment entre cinq cents watts et deux kilowatts par unité~; l'électricité représente, comme noté plus haut, moins de deux pour cent des budgets d'exploitation. Ce qui est crucial pour l'argument de la thèse n'est pas la consommation absolue, mais la structure du coût~: ce coût est discret (lié à la durée d'utilisation), pas continu. Quand les tabulatrices sont éteintes, elles ne consomment pas~; quand les cartes perforées sont dans leurs boîtes, elles n'exigent rien. STORE reste passif, COMPUTE est actif seulement pendant le calcul. La dissémination de l'ICT dans les organisations des années~1930 ne crée pas de consommation continue~: elle crée seulement de l'usage.


\subsection{Boole, Shannon et la logique devenue circuit}\label{subsec:shannon}

\dimmark{D10}{+}%
George Boole publie en 1847 ses \emph{Mathematical Analysis of Logic}, dans lequel il montre que les opérations logiques (et, ou, non, implication) peuvent être représentées par des équations algébriques dont les seules valeurs sont 0 et~1. Le geste est purement mathématique~: Boole ne pense pas au circuit électrique. Il faut attendre quatre-vingt-onze ans pour que Claude Shannon, alors étudiant au MIT, réalise dans sa thèse de master (1938) que les circuits de commutation électrique~--- interrupteurs ouverts ou fermés, courant passant ou non~--- réalisent physiquement les opérations de l'algèbre de Boole. Le~1 logique est le circuit fermé~; le~0 logique est le circuit ouvert. Cette identification transforme l'algèbre de Boole d'un outil mathématique en un programme d'ingénierie~: toute opération logique peut être câblée, et toute suite d'opérations logiques peut être automatisée par des circuits électriques convenablement agencés. Vannevar Bush, directeur du MIT, avait construit dès 1930 un analyseur différentiel électromécanique~; c'est dans ce contexte que Shannon réalise son travail, sous sa direction.

\dimmark{D10}{+}%
Dix ans plus tard, Shannon publie \emph{A Mathematical Theory of Communication} (1948), qui accomplit un second geste~: définir l'information indépendamment de son support physique et de son contenu sémantique. L'information se mesure en bits~; l'entropie informationnelle est formellement analogue à l'entropie thermodynamique de Boltzmann, dont Shannon reprend le terme. Ce second geste fonde la séparation analytique entre l'information et son support~--- séparation qui rendra possible la conception de réseaux universels, indépendants du type de contenu qu'ils transportent, et qui contribuera en même temps à occulter la matérialité et le coût énergétique du signal. La convergence numérique des années~1980--2000 hérite des deux gestes de Shannon~: l'universalité du bit (tout contenu est encodable) et la neutralité du substrat (tout substrat peut transporter n'importe quel bit). Ce que Shannon n'a pas résolu~--- et que \textcite{landauer_irreversibility_1961} tranchera en 1961~--- est précisément ce que cette neutralité cache~: transmettre, stocker et calculer des bits a un coût physique irréductible, que le formalisme de la théorie de l'information rend invisible par construction.

\dimmark{D10}{+}%
Trois ans avant ce papier fondateur, \textcite{wiener_cybernetics_1948} avait publié \emph{Cybernetics, or Control and Communication in the Animal and the Machine}, qui articule l'entropie shannonienne au contrôle par rétroaction et nomme \emph{cybernétique} le champ qui en résulte. Douze ans plus tard, dans \emph{Science}, Wiener met en garde contre la subordination de l'humain à la machine~: les automates qu'il avait contribué à fonder pouvaient, écrit-il, imposer leur logique aux humains qui les emploient si l'on n'y prenait pas garde~\parencite{wiener_consequences_1960}. Le geste qui retournera cette mise en garde en programme d'innovation intervient dès les premières années de la décennie suivante~: \textcite{minsky_steps_1961}, dans \emph{Proceedings of the IRE}, formule la promesse que les machines pourraient, dans le temps d'une vie humaine, surpasser l'humain en intelligence générale. La trajectoire qui mène de l'inquiétude wienerienne à la promesse minskyenne, puis aux investissements massifs des décennies suivantes dans la recherche en intelligence artificielle, illustre une dynamique récurrente dans le déploiement des techniques de l'information~: les imaginaires des fondateurs sont retournés par les acteurs qui les suivent en arguments d'investissement, indépendamment de l'intention initiale. \textcite{johnstone_deep_2026} ont récemment identifié cette dynamique comme structurelle dans la trajectoire de la digitalisation~; le présent chapitre la retrouvera à plusieurs reprises dans les sections qui suivent.


\subsection{Enigma, Bletchley Park et Colossus}\label{subsec:enigma}

\dimmark{D5/D8}{+}%
La Seconde Guerre mondiale accomplit dans le calcul ce que la Première avait accompli dans l'aviation~: elle mobilise des ressources sans limites ordinaires et produit en quelques années des sauts technologiques que le marché aurait mis des décennies à réaliser. En~1939, la Wehrmacht chiffre ses communications avec la machine Enigma, dont le nombre de configurations possibles excède $10^{23}$~; briser Enigma à la main est hors de portée humaine. À Bletchley Park, dix mille personnes~--- mathématiciens, linguistes, ingénieurs, dont Alan Turing~--- travaillent au déchiffrement. La \og~bombe~\fg{} électromécanique que Turing finalise en 1940, fondée sur la logique des contradictions plutôt que sur la force brute, réduit l'espace des configurations à explorer à un volume tractable\footnote{La \og~bombe~\fg{} de Bletchley Park dérive elle-même d'une machine polonaise, la \emph{bomba kryptologiczna}, conçue par Marian Rejewski, Jerzy Różyki et Henryk Zygalski à partir de~1938. Le transfert de technologie de Varsovie à Londres, réalisé quelques semaines avant l'invasion allemande de la Pologne, est l'un des transferts technologiques les plus conséquents de la Seconde Guerre mondiale~; il reste peu connu en dehors de l'histoire des services de renseignement.}. Colossus (1943), conçu par Tommy Flowers de la British Post Office pour décrypter le code Lorenz utilisé par le haut commandement allemand, franchit une étape supplémentaire~: dix tubes à vide, cinq cent brins de papébande perforée lus à cinq mille caractères par seconde, et le premier ordinateur électronique programmable de grande taille qui fonctionne en conditions opérationnelles.

\dimmark{D8/0-TER}{+}%
Colossus fonctionne en continu~: les opérations de déchiffrement ne s'arrêtent pas la nuit, parce que les communications ennemies non plus. C'est la première occurrence dans ce parcours historique d'un système de calcul conçu pour une disponibilité permanente, non par choix architectural mais par nécessité opérationnelle. Le pattern se reproduira~: SAGE tournera sans interruption parce qu'une attaque nucléaire n'attend pas les heures ouvrables (section~\ref{sec:computing})~; les serveurs web fonctionneront sans interruption parce qu'un client potentiel se connecte à n'importe quelle heure (section~\ref{sec:convergence}). La disponibilité permanente, qui produira le couplage énergétique continu décrit dans la section suivante, s'esquisse ici dans le contexte militaire où elle est la moins surprenante~--- et la moins comptabilisée.

\dimmark{D10}{+}%
Bletchley Park illustre enfin une dynamique de financement qui traversera toute l'histoire du calcul~: c'est l'État militaire qui crée le marché technologique que le secteur privé exploitera ensuite, une fois le risque socialisé et les preuves de concept établies. Le secret militaire maintient les résultats hors du domaine public pendant plusieurs décennies~; les ingénieurs de Bletchley Park ne peuvent pas breveter leurs découvertes. Les bénéficiaires du transfert sont les grandes firmes d'électronique~--- IBM, Ferranti, GEC~--- qui recruteront ces ingénieurs après la guerre et s'appuieront sur des savoirs dont le coût de développement a été supporté par le contribuable. Ce mécanisme, que \textcite{mazzucato_entrepreneurial_2013} a systématisé pour la période contemporaine, est structurel dans l'histoire du calcul bien avant qu'elle ne le nomme.


\subsection{Bilan~: la séquence d'électrification des trois fonctions}\label{subsec:bilan_electrification}

\dimmark{0-TER}{*}%
Le parcours de cette section a fait apparaître une séquence qui n'avait pas été nommée jusque-là. Les trois fonctions informationnelles~--- transmettre, stocker, calculer~--- ne s'électrifient pas au même moment~:

\begin{itemize}
\item \textbf{TRANSMIT} s'électrifie dès l'origine~: le télégraphe électrique (1837) \emph{est} électrique par définition. La pile voltaïque, puis le réseau à courant continu et alternatif, portent le signal dès le premier jour. L'électrification de la transmission n'est pas une transition~; c'est une condition initiale.

\item \textbf{COMPUTE} reste mécanique ou manuel pendant plus d'un siècle. La règle à calcul, la calculatrice de Thomas de Colmar, les tabulatrices Hollerith ne tirent pas leur énergie du réseau électrique~: elles sont actionnées à la main ou par des moteurs électriques dont la consommation est marginale dans le bilan d'exploitation. L'électronique de puissance n'entre dans le calcul qu'avec l'ENIAC (1946), soit cent neuf ans après Morse.

\item \textbf{STORE} bascule en dernier. Le papier, le carton perforé, la cire du phonographe, le film photographique, le disque de gomme-laque stockent l'information sans énergie~: l'inscription est une déformation physique permanente du support, et la lecture en est la traduction mécanique ou optique. La mémoire DRAM (Intel 1103, 1970) est la première technologie de stockage qui \emph{exige} un flux d'énergie continu~--- plusieurs centaines de rafraîchissements par seconde~--- pour maintenir son contenu. Avant elle, stocker était gratuit entre l'écriture et la lecture~; après elle, stocker consomme en permanence.
\end{itemize}

\dimmark{0-TER}{*}%
Cette séquence~--- T~(1837) $\rightarrow$ C~(1946) $\rightarrow$ S~(1970)~--- a une conséquence directe pour l'argument de la thèse. Tant que STORE reste passif, le coût énergétique de l'information est ponctuel~: il se produit à la transmission et au calcul, pas entre les deux. L'information peut dormir sur une étagère sans consommer de watt. Le régime énergétique est additif~: chaque usage ajoute une dépense discrète, mais il n'y a pas de fond continu. C'est ce régime que les sections §\ref{sec:telegraph} à §\ref{sec:wireless} ont traversé sans avoir besoin de le nommer, parce qu'il ne posait pas de problème~: l'énergie de l'information était invisible parce qu'elle était négligeable.

\dimmark{0-TER}{+}%
Le basculement vers un régime permanent s'opère en deux temps. Le premier temps est la DRAM (1970)~: la mémoire de travail des ordinateurs devient un flux continu de rafraîchissement. Le second temps est le serveur web (1991) et le datacenter (années~2000)~: l'exigence de disponibilité permanente (\emph{always-on}) transforme STORE d'une ressource ponctuelle en une infrastructure continue. Un fichier stocké sur un serveur web n'existe que si le serveur est allumé~; une page web accessible à tout moment exige une alimentation à tout moment. La section~\ref{sec:convergence} documentera ce basculement~; la section~\ref{sec:contemporain} en mesurera les conséquences. Ce que la présente section établit est le point de départ~: jusqu'en 1970, le substrat électrique des techniques de l'information est une condition d'existence, pas un coût d'exploitation continu. L'occultation de la dimension 0-TER dans les cas du XIX\textsuperscript{e}~siècle n'est pas une erreur d'observation~: elle reflète un régime physique où l'énergie de l'information était réellement négligeable.

\questionch{3}{Comment quantifier la consommation d'énergie de chaque fonction (TRANSMIT, STORE, COMPUTE) à travers les périodes identifiées par ce chapitre~? Le croisement des séries de capacités informationnelles de \textcite{hilbert_worlds_2011} avec les séries d'efficacité énergétique de \textcite{koomey_implications_2011} permettrait-il de reconstituer la consommation agrégée par fonction et par décennie~? Ce croisement n'a jamais été publié. Le chapitre~3 de cette thèse le tentera.}


\begin{table}[htbp]
\caption{Dimensions économiques de l'électrification (§\ref{sec:electrification}).}\label{tab:electrification-dimensions}
\centering\small
\renewcommand{\arraystretch}{1.15}
\begin{tabularx}{\textwidth}{@{}cXc@{}}
\toprule
\textbf{Dim.} & \textbf{Ce que §\ref{sec:electrification} établit} & \textbf{Niveau} \\
\midrule
D2 & Monopole naturel par rendements d'échelle (Insull)~; fragmentation = retard (Londres vs Berlin). & ★ \\
D3 & Capital de réseau irréversible (centrale + distribution)~; intégration verticale AEG. & + \\
D4 & Création d'emplois dans la maintenance et l'exploitation~; déplacement du travail vers la supervision (Nautical Almanac). & + \\
D5 & Demande dérivée (éclairage, puis moteur, puis information)~; diversification des usages pour optimiser le \emph{load factor}. & ★ \\
D6 & Le compteur kilowattheure transforme un flux continu en signal discret facturable~: l'information rend le système de puissance viable. & + \\
D7 & Standardisation des fréquences et tensions comme condition d'interconnexion~; fragmentation londonienne (24 tensions, 10 fréquences). & + \\
D8 & Cycle investissement lourd $\rightarrow$ consolidation (Chicago Edison $\rightarrow$ Commonwealth Edison)~; financement par les banques d'investissement (Deutsche Bank, Drexel Morgan). & ★ \\
D10 & Lock-in DC vs AC (guerre des courants)~; \emph{technological momentum} de Hughes~; \emph{scaling} de Dennard comme bouclier d'occultation. & ★ \\
\midrule
0-TER & Séquence T\,(1837) $\rightarrow$ C\,(1946) $\rightarrow$ S\,(1970)~; régime passif du stockage pré-DRAM~; STORE permanent = consommation continue. & ★ \\
0-GOV & Régulation des monopoles naturels (\emph{public utility commissions})~; divergence Berlin/Londres par la configuration politique, non technique. & ★ \\
\bottomrule
\end{tabularx}
\renewcommand{\arraystretch}{1.0}
\end{table}


% =============================================================================